\titleformat{\chapter}[display]
  {\normalfont\bfseries}{}{0pt}{\Large} % Este comando es para eliminar la aparición de "Capítulo N" en el PDF. Para emplear el formato de la tesis simplemente eliminamos la línea. 
\chapter{Interacciones Intermoleculares}
Las propiedades físico-químicas de los materiales orgánicos empleados en celdas solares dependen en gran medida de las interacciones existentes sobre la morfología de la interfaz, las cuáles, se extienden desde los ángstroms hasta los micrómetros. En escalas moleculares (0.1-5 nm) el empaquetamiento y la orientación individual de cada donador/aceptor controla la formación de grandes agregados que influencian la dinámica del excitón y la movilidad de carga. En escalas nanométricas, (5-50 nm) las interacciones entre donador/aceptor modulan la formación de dominios, incluyendo su tamaño, pureza y el grado de mezcla. A mayores escalas (>50 nm) la morfología es influenciada por una gran segregación de fase y defectos en la interfaz \cite{kilbride2026organic}. \\ \\
Tomando el hilo en la escala molecular, la magnitud de estas interacciones es directamente proporcional a la diferencia entre la energía total de un sistema combinado ($E_{AB}$) y la de sus componentes individuales ($E_A$ y $E_B$), tal como se señala en la Ecuación \ref{ec:int}:
\begin{equation}
\Delta E_{int}=E_{AB}-(E_A+E_B) \label{ec:int} \, .
\end{equation}
La composición de esta energía puede desglosarse en cuatro contribuciones: i) electroestática, ii) intercambio electrónico, iii) inductiva y iv) dispersiva. En la primera, cada monómero tiene una distribución electrónica asociada que le permite producir un campo eléctrico alrededor de ella, la aproximación con otro par adyacente produce una interacción de naturaleza coulómbica, esta contribución es denotada como: $E_{el}$. La segunda componente de la energía de interacción se deriva del requerimiento mecano-cuántico de antisimetría de la función de onda. Al producirse un traslape entre los orbitales de ambos monómeros, el intercambio de electrones entre ellos da lugar al término de energía de intercambio, el cual actúa como una fuerza de repulsión de corto alcance debido al Principio de Exclusión de Pauli. Esta energía es denotada como: $E_{ex}$. En tercer lugar, el acercamiento entre monómeros perturba la densidad electrónica de ambas moléculas produciendo una redistribución de los electrones en el otro. Este efecto es conocido como inducción y puede dividirse en: a) los electrones que pertenecen a \textbf{A} pueden redistribuirse desde su posición original hacia algún espacio desocupado localizado en \textbf{A}, o bien, los electrones de \textbf{A} pueden redistribuirse a un espacio definido por el monómero \textbf{B}, por lo que se produce una transferencia de carga CT de \textbf{A} hacia \textbf{B}. Finalmente, la energía de dispersión surge como una manifestación directa de la correlación electrónica de largo alcance. A diferencia de la inducción, este término se origina cuando uno o más electrones del monómero \textbf{A} se excitan virtualmente hacia orbitales desocupados, un proceso que ocurre de manera acoplada y simultánea con una excitación de energía similar en el monómero \textbf{B}. Estas fluctuaciones correlacionadas de la densidad electrónica generan una fuerza atractiva neta, cuya descripción precisa es indispensable para modelar la estabilidad de los sistemas $\pi$-conjugados \cite{kestner1999basis}. \newpage
\section{Interacciones en Sistemas Aromáticos}
La estructura química de los materiales orgánicos implementados en celdas solares, se encuentra formada principalmente por una serie de anillos conjugados que están fusionados en un sistema ampliamente aromático. En la literatura es común encontrar términos como ''interacciones $\pi-\pi$'' o "$\pi-\pi$ - \textit{stacking}'' haciendo énfasis en un tipo especial de interacción que es único en estos sistemas. No obstante, las interacciones entre moléculas aromáticas de capa cerrada en sus estados electrónicos fundamentales, pueden ser descritas en términos electrostáticos, dispersivos, inductivos, en función del carácter mecanocuántico de repulsión-intercambio y por contribuciones del solvente; EDDIE por sus siglas en inglés \textit{Electrostatic, Dispersion, Desolvation, Induction, and Exchange-repulsion} \cite{xiao2026rethinking}. \\ \\
Las fuerzas intermoleculares en complejos formados por monómeros aromáticos, contribuyen en la formación de tres geometrías que a menudo son discutidas en este tipo de sistemas, veáse la Figura \ref{Dy:2}.
\begin{itemize}
\item \textit{Face-Centered}
\item \textit{T-Shaped}
\item \textit{Parrallel-Displaced}
\end{itemize} 
\begin{figure}[H]
    \centering
    \includegraphics[width=1\textwidth,height=0.3\textwidth]{~/Documentos/Protocolo-Doctorado/imagenes/Dynamics/stacking.png} 
   \caption{Principales dímeros formados en sistemas aromáticos}
    \label{Dy:2}%Esto genera un numero a mi etiqueta
\end{figure} 
Los arreglos de la Figura \ref{Dy:2} son derivados del modelo de Hunter y Sanders \cite{hunter1990nature}, el cuál describe la estructura electrostática de un anillo aromático sencillo, como el benceno, mediante una nube $\pi$ rica en electrones y una periferia deficiente en ellos. Esta distribución genera un marcado momento cuadrupolar (Figura \ref{Dy:1}) \cite{carlone2017physics}, cuya naturaleza electrostática es responsable de que los anillos aromáticos adopten arreglos geométricos específicos cuando los electrones $\pi$ y el fragmento $\sigma$ superan las condiciones desfavorables de la repulsión electrónica \cite{dias2024pi}. Sin embargo, la dispersión es el mayor contribuyente a las interacciones no covalentes entre moléculas aromáticas calculadas en fase gas \cite{xiao2026rethinking}. Por ejemplo, en el benceno se ha reportado experimentalmente por medio de difracción de neutrones, que la forma paralelo-desplazado es la geometría más favorable en distancias menores a 5 $\text{\AA}$, mientras que en separaciones más grandes los anillos aromáticos adyacentes son predominantemente perpendiculares, con dos átomos de H por molécula que se encuentran apuntando hacia los orbitales $\pi$  del otro vecino conjugado \cite{headen2010structure}. \\ \\
En sistemas más grandes se observa un apilamiento centrado en las caras, especialmente en aquellos que presentan fracciones donadoras y aceptoras de electrones en moléculas aromáticas diferentes, a los que a veces se denomina como complejos donador-aceptor. Este comportamiento describe con precisión la naturaleza de los dímeros en la interfaz activa de las celdas solares orgánicas.
\begin{figure}[H]
    \centering
    \includegraphics[width=0.4\textwidth,height=0.3\textwidth]{~/Documentos/Protocolo-Doctorado/imagenes/Dynamics/quadrupole.pdf} 
   \caption{Representación de la distribución de carga del benceno}
    \label{Dy:1}%Esto genera un numero a mi etiqueta
\end{figure} 
\noindent En estos sistemas complejos, el modelo de Interacción Directa Local \textit{Local Direct Interaction} \cite{wheeler2009origin} sugiere que la estabilidad no depende únicamente de la polarización de las densidades $\pi$ aromáticas, sino que son las interacciones electrostáticas directas a través del espacio y las fuerzas dispersivas entre los sustituyentes y los vértices de los anillos adyacentes las que terminan por definir el paisaje energético de la interfaz. No obstante, la parte electrostática dicta la direccionalidad del arreglo, pero la estabilidad energética en fase gaseosa está gobernada predominantemente por la dispersión \cite{xiao2026rethinking}. \\ \\
Por otra parte, los enlaces tradicionales de hidrógeno: II) Figura \ref{Dy:3}; comúnmente denominados como puentes de hidrógeno, constan de interacciones \ce{X-H$\cdots$X, X$=$F, O, N} \cite{haque2023non}. Aunque estas son de naturaleza electrostática, los enlaces muestran restricciones angulares mayores de lo que cabría esperar de interacciones puramente electrostáticas, lo que indica un traslape orbital y un carácter parcialmente covalente \cite{herschlag2018hydrogen}.
\begin{figure}[H]
    \centering
    \includegraphics[width=0.9\textwidth,height=0.55\textwidth]{~/Documentos/Protocolo-Doctorado/imagenes/Dynamics/inter.pdf} 
   \caption{Principales interacciones en dímeros aromáticos}
    \label{Dy:3}%Esto genera un numero a mi etiqueta
\end{figure} 
\noindent El enlace de hidrógeno no tradicional se compone por  arreglos: \ce{O$\cdots$H-C$_{sp^2}$}, \ce{N$\cdots$H-C$_{sp^2}$}, \ce{C-H$\cdots$ $\pi$}, \ce{C-F$\cdots$H-O} y \ce{C-F$\cdots$H-N}. En estudios con el naftaleno y el PBI, estas interacciones son capaces de asistir en la deslocalización de electrones $\pi$, en la planaridad del sistema y en la optimización de la transferencia de carga intramolecular \cite{shang2020perylene}. \\ \\
La incorporación de halógenos (X) o calcógenos (Ch) en las periferias moleculares induce a la formación de un $\sigma$-hole: una región de potencial electrostático positivo localizada en la prolongación del enlace covalente. Esta deficiencia electrónica permite la formación de interacciones altamente direccionales que dictan el autoensamblaje molecular \cite{haque2023non}. El enlace de halógeno actúa como un anclaje molecular hidrofóbico que restringe eficazmente los grados de libertad rotacionales del sistema. Desde una perspectiva fotofísica, esta rigidez estructural es fundamental, ya que suprime las rutas de recombinación no radiativa al limitar los modos vibracionales que disipan la energía. En consecuencia, se observa una prolongación en los tiempos de vida del estado excitado. Además, la presencia de átomos pesados y la formación de estas interacciones pueden potenciar el acoplamiento espín-órbita \cite{yang2020boosting}. \\ \\
En semiconductores que integran unidades de tiofeno, tiazol o EDOT, la formación de enlaces de calcógeno, específicamente contactos de tipo \ce{S \cdots N}, establecen ''anillos virtuales'' de 5 o 6 miembros que bloquean la molécula en conformaciones energéticamente favorables. Dicho control conformacional no solo garantiza la planaridad del esqueleto $\pi$-conjugado, sino que asegura una configuración electrónica óptima para un transporte de carga superior a través de la red cristalina \cite{haque2023non}. Aunado a lo anterior, de las siguientes seis posibles interacciones (\ce{S\cdots N, S\cdots O, S\cdots F, O\cdots N, O\cdots F, N\cdots F}) el tipo  \ce{S \cdots N} mejora de manera significativa el transporte de carga y reduce el $E_g$ \cite{liu2021beyond}. En contraparte, se ha demostrado que la interacción \ce{S\cdots O} es por mucho la interacción más común \cite{haque2023non}. \\ \\
En sistemas derivados de la familia Y-Series, la interacción intramolecular \ce{S \cdots N} en las fracciones de benzo[2,1,3]tiadiazol actúa como un regulador geométrico que favorece un empaquetamiento del tipo \textit{face-to-face}. Esta organización espacial minimiza la formación de niveles trampa derivados del desorden estructural \cite{lin2020non}. En última instancia, la sinergia de estas interacciones no covalentes permite reducir las pérdidas de energía y maximizar la eficiencia cuántica externa (EQE), consolidando a las NCIs como herramientas de diseño indispensables para el desarrollo de celdas solares orgánicas de alto desempeño.
\section{Descripción de Complejos del Sistema BTP-4F-XXXX}
En función de las interacciones no covalentes que hay presentes en sistemas aromáticos, es posible predecir (al menos de manera cualitativa), diversas geometrías entre dímeros donador-aceptor del sistema en estudio. En esta sección se abordará un conjunto de arreglos propuestos entre el sistema de isómeros BTP-4F-T2C8, BTP-4F-T2EH y BTP-4F-T3EH y el donador denotado como D18-Cl. Para simplificar la geometría de estos dímeros se ha optado por representarlos con recuadros sencillos, donde la coloración azul delimita a la fracción que tanto el donador como el aceptor comparten: un núcleo de benzotiadiazol; veáse la Figura \ref{Dy:4}.
\begin{figure}[H]
    \centering
    \includegraphics[width=0.8\textwidth,height=0.7\textwidth]{~/Documentos/Protocolo-Doctorado/imagenes/Dynamics/complex1.png} 
   \caption{Simplificación de los monómeros aceptor y donador}
    \label{Dy:4}%Esto genera un numero a mi etiqueta
\end{figure} 
Dado que tanto el donador D18-Cl como los aceptores isómericos están constituidos por esqueletos conjugados ricos en unidades de tiofeno y heteroátomos de nitrógeno, se favorece la formación de interacciones de calcógeno de tipo \ce{S \cdots N}. Aunado a esto, el grupo ciano en la unidad terminal INIC del aceptor actúa como un potente aceptor de puentes de hidrógeno no tradicionales con las cadenas laterales del D18-Cl. Estas fuerzas sinérgicas son determinantes para estabilizar geometrías de apilamiento $\pi$ específicas, las cuales se han propuesto en cuatro arreglos fundamentales ilustrados en la Figura \ref{Dy:5}. \\ \\
Los dímeros con un arreglo PFF y AFF han sido denotados en función de la orientación relativa de los anillos de benzotiadiazol de ambas moléculas. En el complejo PFF, ambos anillos están en una disposición paralela desplazada. Mientras que en el arreglo AFF este par de benzotiadiazoles apuntan de forma contraria uno respecto al otro. Por otra parte, los dímeros T y L han sido construidos con base a la posición de apilamiento de la unidad terminal del aceptor con respecto al anillo de benzotiadiazol en el complejo T y con respecto al núcleo secundario compuesto por el benzo[1,2-b:4,5-b']ditiofeno o BDT en el arreglo L. 
\begin{figure}[H]
    \centering
    \includegraphics[width=0.8\textwidth,height=0.7\textwidth]{~/Documentos/Protocolo-Doctorado/imagenes/Dynamics/complex2.png} 
   \caption{Propuesta de dímeros donador-aceptor}
    \label{Dy:5}%Esto genera un numero a mi etiqueta
\end{figure} 

%Modelos más actuales como el SAPT, por sus siglas en inglés \textit{Symmetry-Adapted Pertubation Theory} \cite{szalewicz2012symmetry}, y el CCSD(T), \textit{Coupled-Cluster theory} \cite{bartlett2007coupled}, brindan una forma más rigurosa de particionar las energías de interacción no covalentes en términos físicos interpretables. En el SAPT, la energía de interacción es separada en: Electrostática,  Intercambio, Inducción y Dispersión. Todo esto corresponde con el modelo EDDIE, no obstante, la desolvatación es una contribución a la energía libre asociada a la reorganización del disolvente tras la formación del complejo; por lo general, no se incluye en los análisis SAPT en fase gaseosa, pero resulta esencial en contextos de fase condensada e incluso biológicos \cite{xiao2026rethinking}. 